% This is text associated with the open textbook "Open Electromagnetics"
% See immediately below for author, date
% This work is licensed under the Creative Commons Attribution-ShareAlike 4.0 International License. (CC BY-SA 4.0) 
% To view a copy of the license, visit http://creativecommons.org/licenses/by-sa/4.0/

%ORB TITLE Rectangular Waveguide: Propagation Characteristics
%ORB AUTHOR 0        # S.W. Ellingson, Virginia Tech (ellingson@vt.edu)
%ORB DATE 20191202   # yyyymmdd
%ORB ABSTRACT 

\label{m0224_Rectangular_Waveguide_Prop} % <-- Must match name of this file and module directory name.  Don't mess with this!
\index{waveguide!rectangular|(}

In this section, we consider the propagation characteristics of TE and TM modes in rectangular waveguides.
Because these modes exhibit the same phase dependence on $z$, findings of this section apply equally to both sets of modes.
Recall that the TM modes in a rectangular waveguide are given by:
 %
\begin{equation}
\widetilde{E}_z^{(m,n)} = E_0^{(m,n)} \sin\left(\frac{m\pi}{a} x\right) \sin\left(\frac{n\pi}{b} y\right) e^{-jk_z^{(m,n)} z} \label{m0224_eEz}
\end{equation}
%
where 
$E_0^{(m,n)}$ is an arbitrary constant (determined in part by sources), 
and:
% 
\begin{equation}
k_z^{(m,n)} = \sqrt{ \omega^2\mu\epsilon - \left(\frac{m\pi}{a}\right)^2 - \left(\frac{n\pi}{b}\right)^2   } \label{m0224_ekzm}
\end{equation}
%
The TE modes in a rectangular waveguide are:
 %
\begin{equation}
\widetilde{H}_z^{(m,n)} = H_0^{(m,n)} \cos\left(\frac{m\pi}{a} x\right) \cos\left(\frac{n\pi}{b} y\right) e^{-jk_z^{(m,n)} z} \label{m0224_eHz}
\end{equation}
%
where 
$H_0^{(m,n)}$ is an arbitrary constant (determined in part by sources).

{\bf Cutoff frequency}.
First, observe that values of $k_z^{(m,n)}$ obtained from Equation~\ref{m0224_ekzm} are not necessarily real-valued.  
For any given value of $m$, $(k_z^{(m,n)})^2$ will be negative for all values of $n$ greater than some value.
Similarly, for any given value of $n$, $(k_z^{(m,n)})^2$ will be negative for all values of $m$ greater than some value.
Should either of these conditions occur, we find:
% 
\begin{flalign}
\left(k_z^{(m,n)}\right)^2 &=        \omega^2\mu\epsilon - \left(\frac{m\pi}{a}\right)^2 - \left(\frac{n\pi}{b}\right)^2  ~~~ < 0 \nonumber \\
                           &= -\left|\omega^2\mu\epsilon - \left(\frac{m\pi}{a}\right)^2 - \left(\frac{n\pi}{b}\right)^2  \right| \nonumber \\
                           &= -\alpha^2      
\end{flalign}
%
where $\alpha$ is a positive real-valued constant.
So:
% 
\begin{flalign}
k_z^{(m,n)} &= \pm j \alpha
\end{flalign}
%
Subsequently:
%
\begin{flalign}
e^{-jk_z^{(m,n)} z} &= e^{-j \left(\pm j \alpha\right) z} \nonumber \\
                    &= e^{\pm \alpha z} 
\end{flalign}
%
The ``$+$'' sign option corresponds to a wave that grows exponentially in magnitude with increasing $z$, which is non-physical behavior.
Therefore:
%
\begin{equation}
e^{-jk_z^{(m,n)} z} = e^{-\alpha z}
\end{equation}
%
Summarizing:  When values of $m$ or $n$ are such that $(k_z^{(m,n)})^2<0$, the magnitude of the associated wave is no longer constant with $z$.
Instead, the magnitude of the wave decreases exponentially with increasing $z$.
Such a wave does not effectively convey power through the waveguide, and is said to be \emph{cut off}.

Since waveguides are normally intended for the efficient transfer of power, it is important to know the criteria for a mode to be cut off.
Since cutoff occurs when $(k_z^{(m,n)})^2<0$, cutoff occurs when:
%
\begin{equation}
\omega^2\mu\epsilon > \left(\frac{m\pi}{a}\right)^2 + \left(\frac{n\pi}{b}\right)^2
\end{equation}
%
Since $\omega=2\pi f$:
%
\begin{flalign}
f &> \frac{1}{2\pi\sqrt{\mu\epsilon}} \sqrt{ \left(\frac{m\pi}{a}\right)^2 + \left(\frac{n\pi}{b}\right)^2 } \\
  &= \frac{1}{2\sqrt{\mu\epsilon}} \sqrt{ \left(\frac{m}{a}\right)^2 + \left(\frac{n}{b}\right)^2 }
\end{flalign}
%
Note that $1/\sqrt{\mu\epsilon}$ is the phase velocity $v_p$ for the medium used in the waveguide.  With this in mind, let us define:
%
\begin{equation}
v_{pu} \triangleq \frac{1}{\sqrt{\mu\epsilon}}
\end{equation}
%
This is the phase velocity in an unbounded medium having the same permeability and permittivity as the interior of the waveguide.
Thus:
%
\begin{equation}
f > \frac{v_{pu}}{2} \sqrt{ \left(\frac{m}{a}\right)^2 + \left(\frac{n}{b}\right)^2 }
\end{equation}
%
In other words, the mode $(m,n)$ avoids being cut off if the frequency is high enough to meet this criterion.
Thus, it is useful to make the following definition:
%
\begin{equation}
f_{mn} \triangleq \frac{v_{pu}}{2} \sqrt{ \left(\frac{m}{a}\right)^2 + \left(\frac{n}{b}\right)^2 } \label{m0224_efmn}
\end{equation}
%
%%%%%%%%%%%%%%%%%%%%%%%%%%%%%%%%%%%%%%%%%%%%%%%
%ORB HIGHLIGHT
\begin{mdframed}[backgroundcolor=yellow!20,nobreak=true] 
The \emph{cutoff frequency}
\index{cutoff frequency}
$f_{mn}$ (Equation~\ref{m0224_efmn}) is the lowest frequency for which the mode $(m,n)$ is able to propagate (i.e., not cut off). 
\end{mdframed}
%%%%%%%%%%%%%%%%%%%%%%%%%%%%%%%%%%%%%%%%%%%%%%%

%%%%%%%%%%%%%%%%%%%%%%%%%%%%%%%%%%%%%%%%%%%%%%%
%ORB EXAMPLE
~\\ \begin{example} 
\label{m0224_exWR90}
\begin{mdframed}[backgroundcolor=brown!20,linecolor=brown!20]\raggedright
Cutoff frequencies for WR-90. 
\\~\\
WR-90 is a popular implementation of rectangular waveguide.
WR-90 is air-filled with dimensions $a=22.86$~mm and $b=10.16$~mm.
Determine cutoff frequencies and, in particular, the lowest frequency at which WR-90 can be used.

{\bf Solution}. 
Since WR-90 is air-filled, $\mu\approx\mu_0$, $\epsilon\approx\epsilon_0$, and $v_{pu} \approx {1}{\sqrt{\mu_0 \epsilon_0}} \cong 3.00 \times 10^8$~m/s.
Cutoff frequencies are given by Equation~\ref{m0224_efmn}.
Recall that there are no non-zero TE or TM modes with $m=0$ and $n=0$.
Since $a>b$, the lowest non-zero cutoff frequency is achieved when $m=1$ and $n=0$.
In this case, Equation~\ref{m0224_efmn} yields $f_{10} = \underline{6.557~\mbox{GHz}}$; this is the lowest frequency that is able to propagate efficiently in the waveguide.
The next lowest cutoff frequency is $f_{20} = 13.114~\mbox{GHz}$.
The third lowest cutoff frequency is $f_{01} = 14.754~\mbox{GHz}$.
The lowest-order TM mode that is non-zero and not cut off is TM$_{11}$ ($f_{11}=16.145$~\mbox{GHz}).

\end{mdframed}
% note figures will need to go between "\end{mdframed}" and "\end{example}"
\end{example}
%%%%%%%%%%%%%%%%%%%%%%%%%%%%%%%%%%%%%%%%%%%%%%%

{\bf Phase velocity}.
\index{phase velocity}
The phase velocity for a wave propagating within a rectangular waveguide is greater than that of electromagnetic radiation in unbounded space.
For example, the phase velocity of any propagating mode in a vacuum-filled waveguide is greater than $c$, the speed of light in free space.
This is a surprising result.
Let us first derive this result and then attempt to make sense of it.

Phase velocity $v_p$ in the rectangular waveguide is given by
%
\begin{flalign}
v_p &\triangleq \frac{\omega}{k_z^{(m,n)}} \\
    & = \frac{\omega}{\sqrt{ \omega^2\mu\epsilon - \left(m\pi/a\right)^2 - \left(n\pi/b\right)^2   }}
\end{flalign}
%
Immediately we observe that phase velocity seems to be different for different modes.
Dividing the numerator and denominator by $\beta=\omega\sqrt{\mu\epsilon}$, we obtain:
%
\begin{flalign}
v_p &= \frac{1}{\sqrt{\mu\epsilon}}\frac{1}{\sqrt{ 1 - \left(\omega^2\mu\epsilon\right)^{-1}\left[\left(m\pi/a\right)^2 + \left(n\pi/b\right)^2 \right]  }}
\label{m0224_evp2}
\end{flalign}
%
Note that $1/\sqrt{\mu\epsilon}$ is $v_{pu}$, as defined earlier.
Employing Equation~\ref{m0224_efmn} and also noting that $\omega=2\pi f$, Equation~\ref{m0224_evp2} may be rewritten in the following form:
%
\begin{flalign}
v_p &= \frac{v_{pu}}{\sqrt{ 1 - \left(f_{mn}/f\right)^2  }}
\end{flalign}
%
For any propagating mode, $f>f_{mn}$; subsequently, $v_p > v_{pu}$.
In particular, $v_p > c$ for a vacuum-filled waveguide.

How can this not be a violation of fundamental physics?
As noted in Section~\ref{m0176_Phase_and_Group_Velocity}, phase velocity is \emph{not} necessarily the speed at which information travels, but is merely the speed at which a point of constant phase travels.  To send information, we must create a disturbance in the otherwise sinusoidal excitation presumed in the analysis so far.
The complex field structure creates points of constant phase that travel faster than the disturbance is able to convey information, so there is no violation of physical principles.

{\bf Group velocity}.
\index{group velocity}
As noted in Section~\ref{m0176_Phase_and_Group_Velocity}, the speed at which information travels is given by the group velocity $v_g$.
In unbounded space, $v_g=v_p$, so the speed of information is equal to the phase velocity in that case.
In a rectangular waveguide, the situation is different.
We find:
%
\begin{flalign}
v_g &= \left( \frac{\partial k_z^{(m,n)} }{\partial\omega} \right)^{-1} \\
    & = v_{pu}\sqrt{ 1 - \left(f_{mn}/f\right)^2  } \label{m0224_eGV}
\end{flalign} 
%
which is always \emph{less} than $v_{pu}$ for a propagating mode.

Note that group velocity in the waveguide depends on frequency in two ways.
First, because $f_{mn}$ takes on different values for different modes, group velocity is different for different modes.
Specifically, higher-order modes propagate more slowly than lower-order modes having the same frequency.
This is known as \emph{modal dispersion}.
\index{dispersion}
%\index{mode (modal) dispersion}
Secondly, note that the group velocity of any given mode depends on frequency.
This is known as \emph{chromatic dispersion}.
%\index{chromatic dispersion}


%%%%%%%%%%%%%%%%%%%%%%%%%%%%%%%%%%%%%%%%%%%%%%%
%ORB HIGHLIGHT
\begin{mdframed}[backgroundcolor=yellow!20,nobreak=true] 
The speed of a signal within a rectangular waveguide is given by the group velocity of the associated mode (Equation~\ref{m0224_eGV}).
This speed is less than the speed of propagation in unbounded media having the same permittivity and permeability.
Speed depends on the ratio $f_{mn}/f$, and generally decreases with increasing frequency for any given mode.
\end{mdframed}
%%%%%%%%%%%%%%%%%%%%%%%%%%%%%%%%%%%%%%%%%%%%%%%

%%%%%%%%%%%%%%%%%%%%%%%%%%%%%%%%%%%%%%%%%%%%%%%
%ORB EXAMPLE
~\\ \begin{example} 
\begin{mdframed}[backgroundcolor=brown!20,linecolor=brown!20]\raggedright
Speed of propagating in WR-90. 
\\~\\
Revisiting WR-90 waveguide from Example~\ref{m0224_exWR90}:
What is the speed of propagation for a narrowband signal at 10~GHz?

{\bf Solution}. 
Let us assume that ``narrowband'' here means that the bandwidth is negligible relative to the center frequency, so that we need only consider the center frequency.
As previously determined, the lowest-order propagating mode is TE$_{10}$, for which $f_{10} = 6.557~\mbox{GHz}$.
The next-lowest cutoff frequency is $f_{20} = 13.114~\mbox{GHz}$.
Therefore, only the TE$_{10}$ mode is available for this signal.
The group velocity for this mode at the frequency of interest is given by Equation~\ref{m0224_eGV}.
Using this equation, the speed of propagation is found to be \underline{$\cong 2.26 \times 10^8$~m/s}, which is about 75.5\% of $c$.

\end{mdframed}
% note figures will need to go between "\end{mdframed}" and "\end{example}"
\end{example}
%%%%%%%%%%%%%%%%%%%%%%%%%%%%%%%%%%%%%%%%%%%%%%%

\index{waveguide!rectangular|)}


